\section{The exact Schatten-3 threshold}
\label{sec:strip}

\begin{theorem}[Similarity criterion]
\label{thm:similarity}
For every $q\ge1$,
\[
 T_s\in\cS_q
 \quad\Longleftrightarrow\quad
 \sum_{p\in A}p^{-q\Re s}<\infty.
\]
For the prime atom set this is equivalent to $q\Re s>1$.
\end{theorem}

\begin{proof}
The identities
\[
 T_s=ZD_s^AZ^{-1},\qquad D_s^A=Z^{-1}T_sZ
\]
and the two-sided ideal property show that $T_s\in\cS_q$ iff
$D_s^A\in\cS_q$.  The latter is diagonal with singular values
$p^{-\Re s}$.  For primes, the corresponding series converges above
exponent one, diverges at exponent one by Euler's prime harmonic
theorem, and also diverges below it.
\end{proof}

\begin{theorem}[Common reflected strip]
\label{thm:strip}
For every $q\ge1$,
\[
 \cB_s\in\cS_q
 \quad\Longleftrightarrow\quad
 \frac1q<\Re s<1-\frac1q.
\]
The strip is nonempty exactly for $q>2$.  In particular,
\[
 \cB_{1/2+it}\in\bigcap_{q>2}\cS_q,
 \qquad
 \cB_{1/2+it}\notin\cS_2.
\]
\end{theorem}

\begin{proof}
For $B=\bigl(\begin{smallmatrix}0&A\\C&0\end{smallmatrix}\bigr)$,
$B^*B=\diag(C^*C,A^*A)$, so $B\in\cS_q$ iff both legs belong to
$\cS_q$.  \Cref{thm:similarity} gives
$q\Re s>1$ for the first leg and $q(1-\Re s)>1$ for the reflected
leg.  Solving yields the stated strip.
\end{proof}

\begin{proposition}[Critical-line self-adjointness]
\label{prop:selfadjoint}
For every real $t$, $\cB_{1/2+it}$ is compact self-adjoint.
\end{proposition}

\begin{proof}
The incidence idempotents are source-real, and
$1-(1/2+it)=1/2-it=\overline{1/2+it}$.  Therefore
$T_{1-s}^{\sharpop}=T_s^*$ on the line.  Compactness follows from
\cref{thm:strip}, for example with $q=3$.
\end{proof}

\begin{figure}[t]
\centering
\begin{tikzpicture}[
  x=9.8cm,y=1.05cm,
  tick/.style={font=\scriptsize,anchor=north},
  rowlabel/.style={font=\small\bfseries,anchor=east},
  honest/.style={line width=4pt,passgreen},
  excluded/.style={line width=2pt,warningamber,dashed},
  crit/.style={stopred,line width=1pt,densely dashed}
]
  \draw[crit] (0.5,-0.25) -- (0.5,3.3);
  \node[font=\scriptsize,stopred,anchor=south] at (0.5,3.3)
    {$\sigma=\tfrac12$};
  \foreach \y/\q in {2.55/2,1.55/3,0.55/4}{
    \draw[charcoal] (0,\y) -- (1,\y);
    \node[rowlabel] at (-0.04,\y) {$q=\q$};
  }
  \foreach \x/\lab in {0/0,0.25/{1/4},0.333/{1/3},0.5/{1/2},
    0.667/{2/3},0.75/{3/4},1/1}{
    \draw[charcoal] (\x,0.47)--(\x,0.63);
    \node[tick] at (\x,0.43) {$\lab$};
  }
  \draw[excluded] (0.5,2.55)--(0.5,2.55);
  \filldraw[fill=white,draw=warningamber,line width=1.2pt]
    (0.5,2.55) circle (1.8pt);
  \node[font=\scriptsize,warningamber,anchor=west] at (0.53,2.55)
    {empty open strip};
  \draw[honest] (0.333,1.55)--(0.667,1.55);
  \filldraw[fill=white,draw=passgreen] (0.333,1.55) circle (1.6pt);
  \filldraw[fill=white,draw=passgreen] (0.667,1.55) circle (1.6pt);
  \draw[honest] (0.25,0.55)--(0.75,0.55);
  \filldraw[fill=white,draw=passgreen] (0.25,0.55) circle (1.6pt);
  \filldraw[fill=white,draw=passgreen] (0.75,0.55) circle (1.6pt);
  \node[font=\small,formalblue,anchor=west] at (0.68,1.88)
    {first integer-order strip};
\end{tikzpicture}

\caption{The common domain is
$1/q<\sigma<1-1/q$.  At $q=2$ the endpoints meet at the excluded
critical point; $q=3$ is the first integer order giving an open strip
through $\sigma=1/2$.}
\label{fig:strip}
\end{figure}

\Cref{fig:strip} makes the boundary obstruction explicit.  The
critical-line family is not Hilbert--Schmidt, so its second power is
not trace class.  This fact determines which spectral moment and which
modified determinant are admissible.

\begin{remark}[Family versus fixed operator]
\label{rem:family}
\Cref{prop:selfadjoint} does not construct a single operator with $t$
as spectral variable.  It constructs a map
$t\mapsto\cB_{1/2+it}$ whose value is self-adjoint for each $t$.
This is a positive analytic statement but fails the strict fixed-
operator route gate.
\end{remark}
